\documentclass[11pt,a4paper]{article}
\usepackage[utf8]{inputenc}
\usepackage[english]{babel}
\usepackage{amsmath}
\usepackage{amsfonts}
\usepackage{booktabs}
\usepackage{hyperref}
\usepackage{listings}
\usepackage{xcolor}
\usepackage{geometry}

\geometry{margin=1in}

\definecolor{codegreen}{rgb}{0,0.6,0}
\definecolor{codegray}{rgb}{0.5,0.5,0.5}
\definecolor{backcolour}{rgb}{0.95,0.95,0.92}

\lstset{
    backgroundcolor=\color{backcolour},
    commentstyle=\color{codegreen},
    keywordstyle=\color{magenta},
    numberstyle=\tiny\color{codegray},
    stringstyle=\color{purple},
    basicstyle=\ttfamily\small,
    breakatwhitespace=false,
    breaklines=true,
    captionpos=b,
    keepspaces=true,
    numbers=left,
    numbersep=5pt,
    showspaces=false,
    showstringspaces=false,
    showtabs=false,
    tabsize=2,
    frame=single
}

\title{Volumetric Cloud Illumination Report \\ \large Ambient Energy and Multiple Scattering Integration for \texttt{Clouds.hs}}
\author{Advanced Graphics Engine Group}
\date{\today}

\begin{document}

\maketitle

\section{Diagnostic: Why the Clouds Render Too Dark}
The current single-source light model relying strictly on single-scattering Beer's Law ($\text{Transmittance} = e^{-\text{density} \cdot \sigma}$) forces premature light extinction. In real world physics, clouds do not just absorb light; they are highly scattering media with an albedo close to $0.99$. 

When photon units strike a cloud mass, they bounce inward multiple times, illuminating the internal structural cavities. Single-scattering models treat this bouncing energy as lost, causing thick cloud bellies to immediately render jet-black. To resolve this, you must implement two decoupled ambient illumination systems: **Directional Ambient Scattering** and a **Multiple Scattering Octave Approximation**.

---

\section{The Production Solution: Dual-Component Ambient Model}

\subsection{1. Height-Graded Ambient / Sky Term}
Instead of adding physical light entities, sample a vertical gradient ambient term based on the current ray sample height $h$ ($0.0$ at bottom, $1.0$ at top). This simulates the ground-reflected light hitting the bottom and the ambient sky light tinting the top.

\begin{equation}
    I_{\text{ambient}}(h) = \text{lerp}(I_{\text{ground\_albedo}}, I_{\text{sky\_color}}, h) \cdot \text{AmbientStrength}
\end{equation}

Integrate this into the master illumination accumulation loop without running any additional shadow steps:
\begin{equation}
    I_{\text{step}} = (I_{\text{sun}} \cdot \text{Phase}(\theta) \cdot \text{ShadowTransmittance}) + (I_{\text{ambient}}(h) \cdot \text{Density})
\end{equation}

\subsection{2. Multiple Scattering Approximation (The Schneider-Vollnov Trick)}
To simulate light bouncing deeply inside the volume without path-tracing, compute an attenuation profile using three distinct cascading attenuation octaves. Each octave represents a deeper layer of scattering where the cloud density feels effectively lower, and the light attenuation constant decreases:

\begin{equation}
    E_{\text{multi}} = \sum_{k=0}^{2} a^k \cdot e^{-\text{density} \cdot \sigma \cdot b^k}
\end{equation}

Where $a = 0.5$ (energy attenuation per bounce) and $b = 0.25$ (density attenuation per bounce). This mathematically forces thin and medium cloud sections to illuminate from within, mimicking real physical cloud translucency.

---

\section{Haskell FIR Light Loop Adaptation}
Below is the mathematical execution layout transformed into production-grade Haskell \texttt{FIR} syntax to augment the lighting pipeline inside the core loop.

\begin{lstlisting}[language=Haskell, caption={Haskell FIR DSL Multiple Scattering Light Integration}]
-- Multi-Scattering and Ambient Illumination implementation block
calculateVolumetricLighting :: Vec3 Float -> Float -> Float -> Vec3 Float -> Program Pipeline (Vec3 Float)
calculateVolumetricLighting currPos density heightNorm sunDir = do
    sunColor  <- use @"sun_color"
    skyColor  <- use @"sky_ambient_color"
    groundCol <- use @"ground_ambient_color"
    
    -- 1. Compute Height-Graded Ambient Base
    let ambientTerm = lerp groundCol skyColor heightNorm * 0.18
    
    -- Execute standard low-step light march toward sun to find base shadow depth
    shadowDepth <- evaluateLightMarchRay currPos sunDir
    
    -- 2. Three-Octave Multiple Scattering Approximation Array
    let d = shadowDepth * density
    let octave0 = exp (-d * 1.0)
    let octave1 = exp (-d * 0.25) * 0.5
    let octave2 = exp (-d * 0.05) * 0.25
    let totalTransmittance = octave0 + octave1 + octave2
    
    -- 3. Phase functions for forward scattering edge-glow
    let phaseTerm = computeHenyeyGreensteinPhase 0.82
    
    -- Combine Direct Sun Energy and Scattered Atmospheric Ambient
    let directLight  = sunColor * phaseTerm * totalTransmittance
    let indirectLight = ambientTerm * density
    
    return (directLight + indirectLight)
\end{lstlisting}

---

\section{Day-Night Cycle Considerations}
Because your rendering pipeline implements a dynamic day-night cycle, your script parameters must update uniform buffers globally over time:
\begin{itemize}
    \item \textbf{Golden Hour / Sunset:} As the sun passes below the horizon plane, the direct \texttt{sunColor} vector must approach \texttt{Vec3 0.0 0.0 0.0}, while the ambient terms (\texttt{skyColor}) shift from bright cyan to deep twilight navy blue.
    \item \textbf{Night Phase:} The moon can replace the sun vector configuration entirely using a lower global intensity scaling factor ($\approx 0.05$), allowing the clouds to be lit softly by moonlight.
\end{itemize}

---

\section{Feature Evaluation Matrix}

\begin{table}[h!]
\centering
\small
\begin{tabular}{llll}
\toprule
\textbf{Technique} & \textbf{GPU Performance Cost} & \textbf{Visual Realism Quality} & \textbf{Implementation Risk} \\ \midrule
\textbf{More Light Entities} & High ($\mathcal{O}(N \times M)$ loops) & Poor (Looks like plastic) & Disrupts pipeline clusters \\
\textbf{Flat Base Ambient} & Negligible ($\mathcal{O}(1)$ ALU) & Medium (Flat shapes) & Low (Lacks cloud depth) \\
\textbf{Height-Graded Ambient} & Very Low ($\mathcal{O}(1)$ ALU) & High (Anchors clouds to sky) & Low (Highly recommended) \\
\textbf{Multiple Scattering Octaves} & Low (3 additional ALU exp operations) & Maximum (Photorealistic glow) & Medium (Needs exact tuning) \\ \bottomrule
\end{tabular}
\caption{Illumination technique tradeoff configuration analysis matrix.}
\end{table}

\end{document}

